\section{Singular acquired mass and a memory transition}\label{sec:singular}
The following result is an extension of the retained mass-constrained profile, not an assertion that all singular supports have a dimension. It keeps the physically acquired weight of a continuous component through the transition to atomic selection.

\begin{theorem}[Atomic--regular transition]\label{thm:singular}
Consider one regenerative action with unit cost and actual one-call law
\begin{equation}\label{eq:mixture}
 \nu_b=(1-b)\delta_{x_0}+b\nu,\qquad 0<b\le1,
\end{equation}
where $\nu$ is a nonatomic compact $s$-regular probability satisfying \eqref{eq:regular}; $x_0$ is in a fixed bounded task domain. Put
\[
 r=r(N,b)=[1-N(1-b)]_+.
\]
For $N\ge1$ and $m\ge2$,
\begin{equation}\label{eq:singular-rate}
 c\frac{r^{1+2/s}}{(Nb)^{2/s}m^{2/s}}
 \le\Psi_m^{\nu_b}(N)
 \le C\frac{r^{1+2/s}}{(Nb)^{2/s}(m-1)^{2/s}}.
\end{equation}
The constants are independent of $N,m,b$. For any finite $m\ge1$, zero excess is possible exactly when $N(1-b)\ge1$. In particular the complete autonomous risk is $B+\Psi_{M-1}^{\nu_b}(N)$, with no discarded rejected calls.

For $x_0=0$, $\nu$ uniform on $[1,2]$, $0<b\le1/8$, and $1\le N<1/(1-b)$, the single-output profile is exactly
\begin{equation}\label{eq:one-center}
 \Psi_1^{\nu_b}(N)
 =r\left(1+\ell+\frac{\ell^2}{3}\right)
   -r^2\left(1+\frac\ell2\right)^2,
 \qquad \ell=\frac r{Nb}.
\end{equation}
For this uniform law, \eqref{eq:singular-rate} holds with lower constant $1/12$ and upper constant $1/12$ as written with $m$ and $m-1$, respectively.
\end{theorem}
\begin{proof}
Every admissible terminal law has atomic mass at most $N(1-b)$; its continuous mass $u$ is therefore at least $r$. If $u>0$, its normalized continuous part has density at most $Nb/u$ with respect to $\nu$. The regular profile estimate from the proof of Corollary~\ref{cor:regular} gives a distortion at least
\[
 u\,c\left(\frac{Nb}{u}m\right)^{-2/s}
 \ge c\frac{r^{1+2/s}}{(Nb)^{2/s}m^{2/s}}.
\]
The argument applies to arbitrary centers, including a center at the atom. For the upper when $r>0$, put mass $1-r=N(1-b)$ at the atom, use one center there, and use an optimal $(m-1)$-center selected regular law of total mass $r$ and density cap $Nb$. This is admissible because $Nb/r\ge1$. Its distortion gives the upper. When $r=0$, accept only the atom; the mean cost is $1/(1-b)\le N$. Conversely, zero distortion with finitely many centers forces the terminal probability onto atoms, and the only atom has available mass at most $N(1-b)$. Compact attainment from Lemma~\ref{lem:dual} rules out a zero infimum without a zero minimum. This proves the exact threshold.

For \eqref{eq:one-center}, any feasible selected law has continuous mass at most $Nb<1/7$. Its best center, its mean, is therefore below $2/7<1/2$. At that center the atom is strictly closer than every point of $[1,2]$. The elementary bounded-density threshold argument then fills the atom to its capacity and chooses the closest remaining interval, $[1,1+\ell]$, with density $Nb$. To spell out that argument, for a fixed center and threshold $\zeta$, the integral of $(D_C-\zeta)(\widetilde h-h)$ is nonnegative whenever $h$ fills the strict sublevel and partially fills its boundary; equal mass cancels $\zeta$. Thus a jointly optimal center and density has exactly this form. Its continuous mass is $r$, its first moment is $r(1+\ell/2)$, and its second moment is $r(1+\ell+\ell^2/3)$. Subtracting the squared first moment proves the formula.

For the asserted explicit constants, a measure on the line with mass $u$ and Lebesgue density at most $H$ has, for any $m$ centers, distortion at least $u^3/(12H^2m^2)$. Indeed, split it into nearest-center cells of masses $u_j$; filling the interval of length $u_j/H$ symmetrically about its center minimizes each moment and gives $u_j^3/(12H^2)$. Jensen's inequality gives $\sum u_j^3\ge u^3/m^2$. Here $H=Nb$ and $u\ge r$. For the upper put $m-1$ disjoint intervals of length $r/[Nb(m-1)]$ inside $[1,2]$, centered at their midpoints, and fill them at density $Nb$; their total length is at most one. Add the atomic mass. This attains the stated upper constant.
\end{proof}

For a fixed positive small $b$, as $N\uparrow1/(1-b)$, \eqref{eq:one-center} is of order $r$, whereas two or more output labels allow order $r^3$ in the one-dimensional example. These are two distinct finite-memory regimes for the same physical support. The factor $b$ and the deficit $r$ are actual acquisition masses, not a renormalized conditional experiment. The theorem is valid through the zero-risk threshold; it is not obtained by discarding the vanishing stratum.

\begin{proposition}[Action-dependent exact zero threshold]\label{prop:zero}
For the regenerative action family and arbitrary positive scalar costs, $\Psi_m(N;c)=0$ if and only if there are at most $m$ points $C$ and an action $a$ with
\begin{equation}\label{eq:zero-action}
 N\,\frac{\nu_a(C)}{c_a}\ge1.
\end{equation}
Thus mixtures of actions cannot improve the zero-distortion acquisition threshold, although they may improve a positive-distortion tradeoff.
\end{proposition}
\begin{proof}
At a zero minimum, $h_a\nu_a$ is supported on $C$, and
\[
 1\le\sum_a x_a\nu_a(C)
 \le\left(\max_a\frac{\nu_a(C)}{c_a}\right)\sum_a c_ax_a
 \le N\max_a\frac{\nu_a(C)}{c_a}.
\]
One action achieves the finite maximum. Conversely repeatedly use such an action and accept only $C$. The expected cost is $c_a/\nu_a(C)\le N$, and one live label and the $m$ output labels implement it. Continuous parts contribute no mass to $C$ and have not been counted as atoms.
\end{proof}
